\chapter{Status Register}
\label{ch:status-register}

\section{Overview}

The Status Register (\reg{sr}) is a 16-bit special-purpose register that contains condition flags and control bits for
the CPU. It is divided into two bytes with different privilege levels:

\begin{itemize}
	\item \textbf{Low Byte (bits 7--0):} Condition flags (bits 3--0) and reserved bits (bits 7--4) -- readable and directly writable in both modes, and optionally updated as a side effect of ALU instructions in both modes
	\item \textbf{High Byte (bits 15--8):} Control and mode flags -- accessible only in supervisor mode; in protected mode it reads as zero and direct writes to it are silently ignored, with no fault
\end{itemize}

\section{Status Register Layout}

\begin{figure}[H]
	\centering
	\begin{bytefield}[bitwidth=1.5em]{16}
		\bitheader{0-15} \\
		\bitbox{1}{T} &
		\bitbox{1}{A} &
		\bitbox{1}{EA} &
		\bitbox{1}{E4} &
		\bitbox{1}{E3} &
		\bitbox{1}{E2} &
		\bitbox{1}{E1} &
		\bitbox{1}{P} &
		\bitbox{1}{} &
		\bitbox{1}{} &
		\bitbox{1}{} &
		\bitbox{1}{} &
		\bitbox{1}{V} &
		\bitbox{1}{C} &
		\bitbox{1}{N} &
		\bitbox{1}{Z} \\
		\bitbox{8}{\small Privileged} &
		\bitbox{8}{\small Non-Privileged}
	\end{bytefield}
	\caption{Status Register Bit Layout}
	\label{fig:sr-layout}
\end{figure}

\section{Condition Flags (Non-Privileged)}

These flags reflect the results of ALU operations and comparisons. They can be read at any privilege level and are
updated by ordinary instruction execution at any privilege level. A direct write to \reg{sr} can also set them at any
privilege level (see Section~\ref{sec:sr-manual-updates}).

\subsection{Zero Flag (Z) -- Bit 0}

\begin{description}
	\item[Set when:] The result of an operation is zero
	\item[Cleared when:] The result of an operation is non-zero
	\item[Tested by:] Equal (\texttt{==}), Not Equal (\texttt{!=}) condition codes
\end{description}

\subsection{Negative Flag (N) -- Bit 1}

\begin{description}
	\item[Set when:] The result of an operation is negative (bit 15 = 1) when interpreted as a signed two's complement value
	\item[Cleared when:] The result is positive or zero (bit 15 = 0)
	\item[Tested by:] Negative Set (\texttt{NS}), Negative Clear (\texttt{NC}), and signed comparison condition codes
\end{description}

\subsection{Carry Flag (C) -- Bit 2}

\begin{description}
	\item[Set when:] An addition produces a carry out of bit 15, or a subtraction requires a borrow
	\item[Cleared when:] No carry/borrow occurs
	\item[Tested by:] Carry Set (\texttt{CS}), Carry Clear (\texttt{CC}), Unsigned Higher (\texttt{HI}), Unsigned Lower or Same
	      (\texttt{LO}) condition codes
	\item[Note:] For unsigned comparisons, the carry flag acts as an unsigned overflow/borrow indicator
\end{description}

\subsection{Overflow Flag (V) -- Bit 3}

\begin{description}
	\item[Set when:] A signed arithmetic operation produces a result that cannot be represented in 16 bits (signed overflow)
	\item[Cleared when:] No signed overflow occurs
	\item[Tested by:] Overflow Set (\texttt{OS}), Overflow Clear (\texttt{OC}), and signed comparison condition codes
	\item[Note:] Overflow occurs when adding two positive numbers yields a negative result, or adding two negative numbers yields
	      a positive result
\end{description}

\subsection{Reserved Flags -- Bits 4--7}

Bits 4 through 7 of the lower byte are reserved for future use; software must not rely on them. They can be
read/written to but the value of these bits after any instruction that updates the condition flags is undefined.

Consequently, it is recommended to test specific bits with instructions such as \mnemonic{TSAR}/\mnemonic{TSAI}, rather
than do a comparison of the whole status register word with something like \mnemonic{CMPR}/\mnemonic{CMPI}.

\section{Control Flags (Privileged)}

These flags control CPU behavior and can only be modified directly in supervisor mode. In protected mode, a direct
write to \reg{sr} leaves these bits unchanged: the write is silently ignored for bits 15--8 and does not raise a fault.

\subsection{Protected Mode (P) -- Bit 8}

\begin{description}
	\item[When Clear (0):] CPU operates in supervisor mode with full privileges
	\item[When Set (1):] CPU operates in protected mode with restricted access to:
	      \begin{itemize}
		      \item High address-register components (\reg{lh}, \reg{ah}, \reg{sh}, \reg{ph})
		      \item High byte of status register (reads as zero; writes are silently ignored)
		      \item Privileged coprocessor operations (command-word operation nibble \texttt{0x8}--\texttt{0xF}, bits 11--8)
		      \item Software exceptions targeting a vector below \texttt{0x60} (see Chapter~\ref{ch:exceptions})
	      \end{itemize}
	\item[Note:] This bit provides basic memory protection by preventing user code from changing memory segments. Performing a
	      restricted operation raises a privilege violation fault, with the exception of writes to the high byte of the status
	      register, which are silently ignored.
\end{description}

\subsection{Hardware Interrupt Enable Bits (E1-E4) -- Bits 9--12}

These four bits individually control whether each of the four maskable hardware interrupt lines is enabled:

\begin{itemize}
	\item \textbf{E1 (Bit 9):} Enable Hardware Interrupt Line 1 (Level 1, lowest maskable priority)
	\item \textbf{E2 (Bit 10):} Enable Hardware Interrupt Line 2 (Level 2)
	\item \textbf{E3 (Bit 11):} Enable Hardware Interrupt Line 3 (Level 3)
	\item \textbf{E4 (Bit 12):} Enable Hardware Interrupt Line 4 (Level 4, highest maskable priority)
\end{itemize}

\textbf{Note:} Hardware Interrupt Line 5 (Level 5) is a Non-Maskable Interrupt (NMI) and is always enabled. It cannot
be disabled via the status register.

When a bit is set to 1, the corresponding hardware interrupt is enabled and can trigger an exception when requested.
When cleared to 0, interrupts on that line are disabled and are not serviced.

\textbf{Important behavior:}
\begin{itemize}
	\item Disabled interrupts are not queued while the enable bit is 0.
	\item A change to an enable bit takes effect at the next instruction boundary: setting it lets an already-active request in
	      straight after the instruction that set it, and clearing it does not discard a request that was already latched
	      (Section~\ref{sec:pending-hardware-interrupts}).
	\item Hardware interrupt lines 1-4 are level-sensitive, not edge-triggered. If an interrupt request remains active, it will
	      be serviced again as soon as the current handler is finished.
	\item Hardware interrupts can only interrupt execution if their priority level is higher than the currently executing
	      exception level
	\item This prevents the same interrupt from re-triggering while its handler is executing
	\item A device requests service by holding its line until its handler services it. A lower-priority request held during a
	      higher-priority handler is serviced as soon as the exception level drops below its level. The line is not latched while
	      its own handler or a higher-priority one is running, so a request withdrawn before then is not seen
	      (Section~\ref{sec:pending-hardware-interrupts})
\end{itemize}

These bits are typically managed by the operating system's interrupt handling code. At reset, all interrupt enable bits
are cleared (interrupts disabled).

Interrupt handling is discussed in detail in Chapter~\ref{ch:exceptions}.

\subsection{Exception Active (EA) -- Bit 13}

\begin{description}
	\item[When Set (1):] The CPU is currently executing an exception handler (\texttt{current\_exception\_level != 0})
	\item[When Clear (0):] The CPU is executing normal code (\texttt{current\_exception\_level == 0})
	\item[Automatically Set:] When any exception is entered (software, hardware, or fault)
	\item[Set by RETE:] \mnemonic{RETE} sets it if the restored exception level is non-zero and clears it when returning to
	      exception level 0
	\item[Purpose:] Allows software to detect whether it is running within an exception handler context
	\item[Typical Use:] OS code can check this bit to determine appropriate behavior, for example avoiding re-enabling interrupts
	      during nested exception handling
\end{description}

\textbf{Note:} This bit reflects the internal \texttt{current\_exception\_level} state. Programs can read this bit but
cannot directly modify it -- it is managed automatically by the exception handling hardware. This includes
\mnemonic{RETE}, which sets \texttt{EA} from the restored exception level rather than from the saved status register, so
editing bit 13 of a saved status register with \mnemonic{ETTR} has no effect.

\subsection{Trap on Address Overflow (A) -- Bit 14}

\begin{description}
	\item[When Set (1):] An effective-address calculation, address-register auto-update, or program-counter advance that wraps
	      the low word raises a segment overflow fault (vector \texttt{0x03}) before the instruction commits
	      (Section~\ref{sec:segment-overflow-fault})
	\item[When Clear (0):] Address overflows wrap around without any side effects
	\item[Purpose:] When set, helps detect pointer errors during development or can be used to increment/decrement the current
	      segment if a data address goes out of bounds of the segment.
	\item[Typical Use:] Enabled during debugging, disabled in production for performance
\end{description}

\subsection{Trace Mode (T) -- Bit 15}

\begin{description}
	\item[When Set (1):] An instruction trace fault (vector \texttt{0x06}) is generated after every instruction, except at
	      exception level 7
	\item[When Clear (0):] Normal execution
	\item[Purpose:] Allows debuggers to single-step through code
	\item[Note:] Used by SIRCIT (SIRC Instruction Tracer) for debugging
\end{description}

\section{Status Register Updates}

\subsection{Automatic Updates}

The condition flags (Z, N, C, V) are automatically updated by most ALU instructions, unless explicitly disabled:

\begin{itemize}
	\item \textbf{Arithmetic Instructions} (\mnemonic{ADDI}, \mnemonic{SUBI}, \mnemonic{ADCI}, \mnemonic{SBCI}): Update all four condition flags
	\item \textbf{Logical Instructions} (\mnemonic{ANDI}, \mnemonic{ORRI}, \mnemonic{XORI}): Update Z and N flags; C and V flags are cleared
	\item \textbf{Compare Instructions} (\mnemonic{CMNI}, \mnemonic{CMNR}, \mnemonic{CNCI}, \mnemonic{CNCR},
	      \mnemonic{CMPI}, \mnemonic{CMPR}, \mnemonic{CPCI}, \mnemonic{CPCR}): Same as Arithmetic Instructions but does not update destination
	\item \textbf{Test Instructions} (\mnemonic{TSAI}, \mnemonic{TSAR}, \mnemonic{TSOI}, \mnemonic{TSOR},
	      \mnemonic{TSXI}, \mnemonic{TSXR}): Same as Logical Instructions but does not update destination
	\item \textbf{Move Instructions} (\mnemonic{MOVI}, \mnemonic{MOVR}): Same as Logical Instructions. They do not accept
	      \texttt{[S]}, since they have no shifted operand; an encoding of the opcode with AF \texttt{10}, not reachable
	      from public assembly syntax, leaves all four flags architecturally undefined (\texttt{U U U U}).
	\item \textbf{Load Instructions} (\mnemonic{LOAD}, \mnemonic{LDEA}): Never update the condition flags and have no
	      status source encoding to override.
	\item \textbf{ALU Instructions with a Shift Operand}: The flags come from the ALU result by default. They come from the
	      shift instead when the \texttt{[S]} status override modifier is used, or with the \mnemonic{SHFT} meta instruction.
\end{itemize}

\subsection{Manual Updates}
\label{sec:sr-manual-updates}

In supervisor mode, the status register can be directly manipulated like any other register:

% qa:assemble
\begin{lstlisting}
; Set protected mode bit (only in supervisor mode)
ORRI sr, #0x0100

; Enable all maskable hardware interrupts (set bits 9-12)
; Note: Hardware interrupt line 5 is an NMI and is always enabled
ORRI sr, #0x1E00

; Disable all maskable hardware interrupts (clear bits 9-12)
ANDI sr, #0xE1FF

; Restore status register saved to stack
LOAD sr, (s)+

\end{lstlisting}

An ALU instruction that targets \reg{sr} writes its result to \reg{sr} and then overwrites bits 7--0 with the flags
produced by that same instruction, so only bits 15--8 can be set this way. To write the low byte, suppress the flag
update with the \texttt{[N]} status source (for example \mnemonic{ORRI[N]} \reg{sr}, \imm{0x0003}) or use
\mnemonic{MOVR[N]} \reg{sr}, \texttt{rN}. A direct write to \reg{sr} also preserves the current value of bit 13
(\texttt{EA}); only exception entry and \mnemonic{RETE} change it, so the \mnemonic{LOAD} \reg{sr}, (s)+ restore above
does not reproduce the exception-active state saved at the time of the original write.

\textbf{Important:} When running in protected mode, a direct write to \reg{sr} updates only the low byte (bits 7--0).
Bits 15--8 -- including the Protected Mode bit itself -- are left unchanged, silently and without a fault, so a
protected-mode program cannot leave protected mode or alter the interrupt-enable, trace, or address-overflow bits this
way. Because bits 15--8 read as zero in protected mode, a read-modify-write of \reg{sr} behaves consistently: the high
byte is never disturbed. This applies to every instruction that names \reg{sr} as its destination, including the
flag-only comparison and test instructions (\mnemonic{CMNI}, \mnemonic{CMNR}, \mnemonic{CNCI}, \mnemonic{CNCR},
\mnemonic{CMPI}, \mnemonic{CMPR}, \mnemonic{CPCI}, \mnemonic{CPCR}, \mnemonic{TSAI}, \mnemonic{TSAR},
\mnemonic{TSOI}, \mnemonic{TSOR}, \mnemonic{TSXI}, \mnemonic{TSXR}), which do not write their destination and are
therefore unaffected.

\section{Reading the Status Register}

In supervisor mode, reading \reg{sr} returns the full 16-bit value. In protected mode, reading \reg{sr} returns the
lower byte in bits 7--0, with bits 15--8 masked to zero (to hide supervisor-only state).

% qa:assemble
\begin{lstlisting}
; In supervisor mode - gets full 16 bits
MOVR r1, sr  ; r1 = 0xXXXX (all bits)

; In protected mode - high byte masked
MOVR r1, sr  ; r1 = 0x00XX (high byte forced to 0)
\end{lstlisting}

\section{Exception Handling Effects}

When an exception occurs:

\begin{enumerate}
	\item The current status register is saved
	\item Protected mode bit (P) is cleared (entering supervisor mode)
	\item Trace mode bit (T) is cleared (disabling single-stepping in exception handler)
	\item Exception active bit (EA) is set
	\item The hardware interrupt enable bits are not changed. A pending interrupt is serviced during a handler only if its
	      priority level is higher than the current exception level
\end{enumerate}

When returning from an exception with \mnemonic{RETE}, the saved status register is restored.

\section{Usage Examples}

The SIRC-1 provides syntax to manually control how status flags are updated during ALU instructions, including those
with a shift operand. This is particularly useful for advanced control flow and flag preservation. See
Section~\ref{sec:status-override} in Chapter~\ref{ch:shift-operations} for complete details and worked examples.

Condition-code-driven branching and conditional execution (comparing registers, testing bits, and acting on the result)
are a property of the condition flags this chapter defines, but the idioms themselves belong to
Chapter~\ref{ch:condition-codes} and the individual instruction entries (for example, \mnemonic{TSAI}/\mnemonic{TSAR}
in Chapter~\ref{ch:alu-instructions}), which cover them in more depth than a status-register-focused example could.

\begin{example}
	\begin{lstlisting}
; Entering protected mode (from supervisor mode)
ORRI sr, #0x0100  ; Set protected mode bit
; Now running in protected mode

; Attempting to write privileged register will fault
; ADDI ah, #0x05  ; This would cause exception!
	\end{lstlisting}
	\caption{Enabling Protected Mode}
\end{example}

\begin{example}
	\begin{lstlisting}
; Save status register (supervisor mode only -- see the manual updates section)
MOVR r7, sr

; Do some operations that change flags
ADDR r1, r2, r3
CMPI r1, #100

; Restore original status ([N] stops the flag update overwriting the low byte)
MOVR[N] sr, r7
	\end{lstlisting}
	\caption{Saving and Restoring Status}
\end{example}

As with the stack restore in Section~\ref{sec:sr-manual-updates}, this restore does not reproduce bit 13 (\texttt{EA});
only exception entry and \mnemonic{RETE} change it.
